\scanpage{135}
\mspage{51}

pour que $M$ admette assez d'injectifs -- ce qui implique alors l'existence d'une petite sous-catégorie strictement \uline{co}génératrice, donc que $M^{\mathrm{o}}$ est également un pseudo-topos.$^{*}$

% slanted
\margin{C'est le cas notamment si $M$ est un \uline{topos}.}

Je \unsure{visualiserai|visualise} les \inserted{catégories} pseudo-topos comme étant la catégorie des ``faisceaux'' sur un objet $\unsure{X|\mathcal X}$, visualisé comme une sorte \unsure{$^{*}$|``} d'``espace'', que je distingue de la catégorie $M=\operatorname{Faisc}(X)$ des faisceaux sur $X$. \deleted{Les} C'est $X$ qui sera pour moi le ``pseudo-topos''. Les pseudo-topos forment une 2-catégorie, la variance des $X$ étant opposée à celle des \deleted{\unsure{faisc}} $\operatorname{Fais}(X)=M$. \deleted{$\underline{\operatorname{Hom}}(X,Y)={}$} De façon précise, si $X$ correspond à $M$, $\unsure{Y|\mathcal Y}$ à $N$, un morphisme $/f/$ de $X$ dans $Y$ sera donné par un foncteur

% slanted
\margin{\unsure{Dire|D'où} ici que les \deleted{\unsure{faisceaux sur}} objets de $M$ sont visualisés comme des \deleted{\unsure{obj.}} ``\unsure{sections|objets}'' généralisés de $X$. On a $M\xrightarrow{\approx}\underline{\operatorname{Hom}}^{!}(M^{\mathrm{o}},\mathrm{Ens})$, ici $M$ s'\unsure{interprète|réinterprète} bel et bien comme les ``faisceaux'' sur cette catégorie $M$ d'\unsure{ensembl.|arguments}\ldots{} cf.}

\[
f^{*}:\deleted{\unsure{M}}N\longrightarrow M
\tag*{(4.1.1)}
\]

en sens inverse, commutant \deleted{\unsure{à}} aux \uline{lim} inductives quelconques. [\uline{NB} si $X,Y$ sont $=$ des topos, i.e. $M,N$ des catégories-topos, \deleted{on exige} et si on veut que $f$ soit un morphisme de topos, on \uline{exige} de plus que $f^{*}$ soit de plus \uline{exact à g.} On dira qu'un \uline{morphisme de pseudo-topos est} \uuline{exact} si de plus $f^{*}$ est exact à g.]
% archive sheet number: 134
